% Source: https://artofproblemsolving.com/community/c3h1783316_diophantine_equation
% Stable identity: aops-topic-1783316-q01
Yes; it has infinitely many. By Bezout's identity there is at least one solution $(x_0,y_0)$. For every integer $t$,
\[
a(x_0+bt)+b(y_0-at)=ax_0+by_0+abt-abt=1,
\]
so $(x_0+bt,\ y_0-at)$ is a solution, and different $t$ give different solutions because $b\ne0$.

These are all the solutions. If $ax+by=1$, subtracting $ax_0+by_0=1$ gives $a(x-x_0)=-b(y-y_0)$. So $b$ divides $a(x-x_0)$, and since $\gcd(a,b)=1$, $b\mid x-x_0$; write $x-x_0=bt$. Then $abt=-b(y-y_0)$, so $y-y_0=-at$.

Answer: $\boxed{\text{yes: the solutions are }x=x_0+bt,\ y=y_0-at\ (t\in\mathbb Z)}$, infinitely many. For example $3x+5y=1$ has the solutions $(2,-1),(7,-4),(-3,2),\dots$.
